% Appendices
\chapter{APPENDICES}

\section{Appendix A: Code Samples}

\subsection{A.1 Workflow Generation Algorithm}

\begin{lstlisting}[language=TypeScript, caption=Workflow Generation Function]
interface WorkflowStep {
  stepOrder: number
  requiredRole: 'faculty' | 'hod' | 'coe' | 'principal'
  description: string
}

function analyzeDocumentContext(document: DocumentMetadata): WorkflowStep[] {
  const steps: WorkflowStep[] = []
  
  // Notices and Circulars
  if (document.type === 'notice' || document.type === 'circular') {
    if (document.scope === 'department') {
      steps.push({ 
        stepOrder: 1, 
        requiredRole: 'hod', 
        description: 'HOD Approval' 
      })
    } else if (document.scope === 'college') {
      steps.push({ 
        stepOrder: 1, 
        requiredRole: 'hod', 
        description: 'HOD Approval' 
      })
      steps.push({ 
        stepOrder: 2, 
        requiredRole: 'principal', 
        description: 'Principal Approval' 
      })
    } else if (document.scope === 'institution') {
      steps.push({ 
        stepOrder: 1, 
        requiredRole: 'hod', 
        description: 'HOD Approval' 
      })
      steps.push({ 
        stepOrder: 2, 
        requiredRole: 'coe', 
        description: 'COE Review' 
      })
      steps.push({ 
        stepOrder: 3, 
        requiredRole: 'principal', 
        description: 'Principal Approval' 
      })
    }
  }
  
  // Policy Documents
  else if (document.type === 'policy') {
    steps.push({ 
      stepOrder: 1, 
      requiredRole: 'hod', 
      description: 'HOD Review' 
    })
    steps.push({ 
      stepOrder: 2, 
      requiredRole: 'principal', 
      description: 'Principal Approval' 
    })
  }
  
  // Exam Schedules
  else if (document.type === 'exam_schedule') {
    steps.push({ 
      stepOrder: 1, 
      requiredRole: 'hod', 
      description: 'HOD Review' 
    })
    steps.push({ 
      stepOrder: 2, 
      requiredRole: 'coe', 
      description: 'COE Verification' 
    })
    steps.push({ 
      stepOrder: 3, 
      requiredRole: 'principal', 
      description: 'Principal Approval' 
    })
  }
  
  return steps
}
\end{lstlisting}

\subsection{A.2 Digital Signature Embedding}

\begin{lstlisting}[language=TypeScript, caption=PDF Signature Embedding]
import { PDFDocument } from 'pdf-lib'

async function embedSignatureInPDF(
  pdfBytes: Uint8Array,
  signatureImageBytes: Uint8Array,
  x: number,
  y: number,
  width: number = 150,
  height: number = 50
): Promise<Uint8Array> {
  try {
    // Load the PDF document
    const pdfDoc = await PDFDocument.load(pdfBytes)
    
    // Embed the signature image (PNG format)
    const signatureImage = await pdfDoc.embedPng(signatureImageBytes)
    
    // Get the first page
    const firstPage = pdfDoc.getPage(0)
    
    // Draw the signature image on the page
    firstPage.drawImage(signatureImage, {
      x: x,
      y: y,
      width: width,
      height: height,
    })
    
    // Save and return the modified PDF
    const modifiedPdfBytes = await pdfDoc.save()
    return modifiedPdfBytes
  } catch (error) {
    console.error('Error embedding signature:', error)
    throw new Error('Failed to embed signature in PDF')
  }
}
\end{lstlisting}

\subsection{A.3 Code Execution with Judge0}

\begin{lstlisting}[language=TypeScript, caption=Code Execution Function]
async function executeCodeWithJudge0(
  sourceCode: string,
  languageId: number,
  testCases: TestCase[]
): Promise<ExecutionResult[]> {
  const results: ExecutionResult[] = []
  
  for (const testCase of testCases) {
    try {
      // Submit code for execution
      const submitResponse = await fetch(
        'https://judge0-ce.p.rapidapi.com/submissions',
        {
          method: 'POST',
          headers: {
            'Content-Type': 'application/json',
            'X-RapidAPI-Key': process.env.JUDGE0_API_KEY!,
            'X-RapidAPI-Host': 'judge0-ce.p.rapidapi.com'
          },
          body: JSON.stringify({
            source_code: sourceCode,
            language_id: languageId,
            stdin: testCase.input,
            expected_output: testCase.expectedOutput
          })
        }
      )
      
      const submission = await submitResponse.json()
      const token = submission.token
      
      // Poll for results
      let status = 'In Queue'
      let result
      
      while (status === 'In Queue' || status === 'Processing') {
        await new Promise(resolve => setTimeout(resolve, 1000))
        
        const resultResponse = await fetch(
          \`https://judge0-ce.p.rapidapi.com/submissions/\${token}\`,
          {
            headers: {
              'X-RapidAPI-Key': process.env.JUDGE0_API_KEY!,
              'X-RapidAPI-Host': 'judge0-ce.p.rapidapi.com'
            }
          }
        )
        
        result = await resultResponse.json()
        status = result.status.description
      }
      
      results.push({
        testCaseId: testCase.id,
        passed: result.status.id === 3, // 3 = Accepted
        output: result.stdout || result.stderr,
        executionTime: result.time,
        memory: result.memory
      })
    } catch (error) {
      results.push({
        testCaseId: testCase.id,
        passed: false,
        error: 'Execution failed'
      })
    }
  }
  
  return results
}
\end{lstlisting}

\section{Appendix B: Database Schema}

\subsection{B.1 Complete Schema SQL}

\begin{lstlisting}[language=SQL, caption=Database Schema Creation, basicstyle=\ttfamily\footnotesize]
-- Users table
CREATE TABLE users (
    id UUID PRIMARY KEY DEFAULT uuid_generate_v4(),
    email TEXT UNIQUE NOT NULL,
    full_name TEXT NOT NULL,
    role TEXT NOT NULL CHECK (role IN (
        'student', 'faculty', 'hod', 'coe', 'principal', 'admin',
        'accountant', 'librarian', 'admission_incharge'
    )),
    department TEXT,
    section TEXT,
    employee_id TEXT,
    phone TEXT,
    current_year INTEGER,
    current_semester INTEGER,
    status TEXT DEFAULT 'pending',
    approval_stage TEXT,
    created_by_email TEXT,
    created_at TIMESTAMP DEFAULT NOW()
);

-- Documents table
CREATE TABLE documents (
    id UUID PRIMARY KEY DEFAULT uuid_generate_v4(),
    title TEXT NOT NULL,
    type TEXT NOT NULL,
    category TEXT,
    department TEXT NOT NULL,
    scope TEXT NOT NULL,
    creator_id UUID REFERENCES users(id),
    file_url TEXT NOT NULL,
    status TEXT DEFAULT 'pending',
    created_at TIMESTAMP DEFAULT NOW(),
    updated_at TIMESTAMP DEFAULT NOW()
);

-- Workflows table
CREATE TABLE workflows (
    id UUID PRIMARY KEY DEFAULT uuid_generate_v4(),
    document_id UUID REFERENCES documents(id) ON DELETE CASCADE,
    policy_valid BOOLEAN DEFAULT TRUE,
    generated_steps JSONB NOT NULL,
    created_at TIMESTAMP DEFAULT NOW()
);

-- Approvals table
CREATE TABLE approvals (
    id UUID PRIMARY KEY DEFAULT uuid_generate_v4(),
    workflow_id UUID REFERENCES workflows(id) ON DELETE CASCADE,
    step_order INTEGER NOT NULL,
    required_role TEXT NOT NULL,
    approver_id UUID REFERENCES users(id),
    status TEXT DEFAULT 'pending',
    signature_url TEXT,
    comment TEXT,
    acted_at TIMESTAMP,
    created_at TIMESTAMP DEFAULT NOW()
);

-- Assignments table
CREATE TABLE assignments (
    id UUID PRIMARY KEY DEFAULT uuid_generate_v4(),
    title TEXT NOT NULL,
    description TEXT,
    type TEXT NOT NULL CHECK (type IN (
        'code', 'link', 'project', 'document'
    )),
    faculty_email TEXT NOT NULL,
    department TEXT,
    section TEXT,
    max_score INTEGER DEFAULT 100,
    deadline TIMESTAMP,
    created_at TIMESTAMP DEFAULT NOW()
);

-- Test cases table
CREATE TABLE test_cases (
    id UUID PRIMARY KEY DEFAULT uuid_generate_v4(),
    assignment_id UUID REFERENCES assignments(id) ON DELETE CASCADE,
    input TEXT NOT NULL,
    expected_output TEXT NOT NULL,
    points INTEGER DEFAULT 10,
    is_hidden BOOLEAN DEFAULT FALSE
);

-- Submissions table
CREATE TABLE submissions (
    id UUID PRIMARY KEY DEFAULT uuid_generate_v4(),
    assignment_id UUID REFERENCES assignments(id) ON DELETE CASCADE,
    student_email TEXT NOT NULL,
    code TEXT,
    language TEXT,
    link_url TEXT,
    file_url TEXT,
    score INTEGER,
    status TEXT DEFAULT 'pending',
    submitted_at TIMESTAMP DEFAULT NOW()
);

-- Create indexes for performance
CREATE INDEX idx_documents_creator ON documents(creator_id);
CREATE INDEX idx_documents_status ON documents(status);
CREATE INDEX idx_approvals_workflow ON approvals(workflow_id);
CREATE INDEX idx_approvals_status ON approvals(status);
CREATE INDEX idx_submissions_assignment ON submissions(assignment_id);
CREATE INDEX idx_submissions_student ON submissions(student_email);
\end{lstlisting}

\section{Appendix C: User Manual}

\subsection{C.1 Student Guide}

\textbf{Logging In:}
\begin{enumerate}
\item Navigate to the application URL
\item Enter your registered email and password
\item Click "Login" button
\end{enumerate}

\textbf{Submitting Assignments:}
\begin{enumerate}
\item Click "Assignments" in the navigation bar
\item Find the assignment you want to submit
\item Click on the assignment title
\item For code assignments:
    \begin{itemize}
    \item Write or paste your code
    \item Select the programming language
    \item Click "Run Code" to test
    \item Review results and click "Submit"
    \end{itemize}
\item For other types, upload file or enter link and click "Submit"
\end{enumerate}

\textbf{Viewing Grades:}
\begin{enumerate}
\item Go to "Assignments" page
\item Your score is displayed next to each submitted assignment
\item Click "View Details" to see test case results
\end{enumerate}

\subsection{C.2 Faculty Guide}

\textbf{Creating Assignments:}
\begin{enumerate}
\item Click "Assignments" → "Create Assignment"
\item Fill in title and description
\item Select assignment type
\item For code assignments, add test cases
\item Set deadline and max score
\item Click "Create Assignment"
\end{enumerate}

\textbf{Grading Submissions:}
\begin{enumerate}
\item Go to assignment detail page
\item Click "View Submissions"
\item For auto-graded assignments, scores are automatic
\item For manual grading, enter score and feedback
\item Click "Save Grade"
\end{enumerate}

\textbf{Uploading Documents:}
\begin{enumerate}
\item Click "Upload Document"
\item Fill in document metadata
\item Select file to upload
\item Click "Submit for Approval"
\item Track status on Dashboard
\end{enumerate}

\subsection{C.3 Approver Guide (HOD/Principal)}

\textbf{Reviewing Documents:}
\begin{enumerate}
\item Pending approvals appear on Dashboard
\item Click on document to review
\item Download and review the file
\item Capture or upload your signature
\item Add optional comments
\item Click "Approve" or "Reject"
\end{enumerate}

\textbf{Managing Department:}
\begin{enumerate}
\item Go to "Departments" page
\item View department structure
\item Add or edit sections
\item Assign users to departments/sections
\end{enumerate}

\section{Appendix D: API Documentation}

\subsection{D.1 Authentication Endpoints}

\textbf{POST /api/auth/login}

Request:
\begin{lstlisting}[language=JSON]
{
  "email": "user@example.com",
  "password": "password123"
}
\end{lstlisting}

Response:
\begin{lstlisting}[language=JSON]
{
  "success": true,
  "user": {
    "id": "uuid",
    "email": "user@example.com",
    "role": "faculty"
  },
  "token": "jwt-token-here"
}
\end{lstlisting}

\subsection{D.2 Document Endpoints}

\textbf{POST /api/documents/upload}

Headers: \texttt{Authorization: Bearer <token>}

Request (multipart/form-data):
\begin{itemize}
\item title: string
\item type: string
\item department: string
\item scope: string
\item file: File
\end{itemize}

Response:
\begin{lstlisting}[language=JSON]
{
  "success": true,
  "documentId": "uuid",
  "workflowId": "uuid"
}
\end{lstlisting}

\subsection{D.3 Assignment Endpoints}

\textbf{POST /api/assignments/create}

Headers: \texttt{Authorization: Bearer <token>}

Request:
\begin{lstlisting}[language=JSON]
{
  "title": "Assignment Title",
  "description": "Description",
  "type": "code",
  "department": "Computer Science",
  "section": "A",
  "deadline": "2024-12-31T23:59:59Z",
  "testCases": [
    {
      "input": "5",
      "expectedOutput": "120",
      "points": 10
    }
  ]
}
\end{lstlisting}

Response:
\begin{lstlisting}[language=JSON]
{
  "success": true,
  "assignmentId": "uuid"
}
\end{lstlisting}

\section{Appendix E: Screenshots}

\textit{Note: Add actual screenshots of your application here}

\begin{enumerate}
\item Login Page
\item Student Dashboard
\item Faculty Dashboard
\item Document Upload Form
\item Approval Interface with Signature Pad
\item Assignment Creation Form
\item Code Submission Interface
\item Notice Board
\item Department Management
\item Analytics Dashboard
\end{enumerate}

\section{Appendix F: System Requirements}

\subsection{F.1 Server Requirements (Production)}

\begin{table}[H]
\centering
\begin{tabular}{|l|l|}
\hline
\textbf{Component} & \textbf{Specification} \\
\hline
Hosting Platform & Vercel (recommended) or equivalent \\
Database & PostgreSQL 15+ (Supabase) \\
Storage & Cloud storage (Supabase/AWS S3) \\
CPU & 2+ vCPUs (auto-scaling) \\
RAM & 4GB+ (auto-scaling) \\
Storage Space & 50GB+ (expandable) \\
Bandwidth & Unlimited (CDN) \\
SSL Certificate & Required (automatic via Let's Encrypt) \\
\hline
\end{tabular}
\caption{Server Requirements}
\end{table}

\subsection{F.2 Client Requirements}

\begin{table}[H]
\centering
\begin{tabular}{|l|l|}
\hline
\textbf{Component} & \textbf{Specification} \\
\hline
Device & Desktop, Laptop, Tablet, or Smartphone \\
Operating System & Windows 10+, macOS 10.15+, iOS 13+, Android 9+ \\
Browser & Chrome 90+, Firefox 88+, Safari 14+, Edge 90+ \\
RAM & 2GB minimum, 4GB recommended \\
Screen Resolution & 360px+ width (mobile responsive) \\
Internet Connection & 1Mbps+ (5Mbps+ recommended) \\
\hline
\end{tabular}
\caption{Client Requirements}
\end{table}

\section{Appendix G: Glossary}

\begin{description}
\item[API] Application Programming Interface - A set of protocols for building software applications
\item[CDN] Content Delivery Network - Distributed server network for faster content delivery
\item[CRUD] Create, Read, Update, Delete - Basic database operations
\item[HOD] Head of Department - Senior faculty member managing a department
\item[JWT] JSON Web Token - Compact token format for secure information exchange
\item[OAuth] Open Authorization - Industry-standard protocol for authorization
\item[PDF] Portable Document Format - File format for document exchange
\item[REST] Representational State Transfer - Architectural style for web services
\item[Row-Level Security] Database feature restricting data access at row level
\item[SSR] Server-Side Rendering - Generating HTML on the server
\item[SSL/TLS] Secure Sockets Layer / Transport Layer Security - Encryption protocols
\item[UUID] Universally Unique Identifier - Standard for generating unique IDs
\end{description}

\section{Appendix H: Acknowledgment of Open Source}

This project uses the following open-source technologies:

\begin{itemize}
\item Next.js (MIT License)
\item React (MIT License)
\item TypeScript (Apache License 2.0)
\item Tailwind CSS (MIT License)
\item pdf-lib (MIT License)
\item PostgreSQL (PostgreSQL License)
\item Supabase (Apache License 2.0)
\end{itemize}

We acknowledge and thank the open-source community for their invaluable contributions.

\newpage
